\uuid{i7Mq}
\exo7id{1193}
\titre{exo7 1193}
\auteur{legall}
\organisation{exo7}
\datecreate{1998-09-01}
\isIndication{false}
\isCorrection{false}
\chapitre{Suite}
\sousChapitre{Convergence}
\cpge{mpsi-mp2i.nombres-reels-suites.d, pcsi.nombres-reels-suites.c, ptsi.nombres-reels-suites.c, mpsi-mp2i.nombres-reels-suites.i, pcsi.nombres-reels-suites.g, ptsi.nombres-reels-suites.g, mpsi-mp2i.calcul-algebrique-trigonometrie.d, pcsi.calcul-algebrique-trigonometrie.d, ptsi.calcul-algebrique-trigonometrie.d, bcpst1.suites-reelles, bcpst1.suites-reelles-usuelles}
\module{Analyse}
\niveau{L1}
\difficulte{}

\contenu{
\texte{
Soient $  (u_n)_{n\geq 2}  $
d\'efinie par $  \displaystyle{ u_n =\prod_{k=2}^n
\hbox{cos} (\frac{\pi}{2^k})}  $ et $  \displaystyle{v_n =u_n \hbox{sin}(\frac{\pi} { 2^n})}  .$
}
\begin{enumerate}
    \item \question{Montrer que $  (u_n)_{n\geq 2}  $ est convergente.}
    \item \question{Montrer que $  (v_n)_{n\geq 2}  $ est une suite g\'eom\'etrique. En d\'eduire
la limite de $  (u_n)_{n\geq 2}  .$}
\end{enumerate}
}
